\documentclass{article}
%PACKAGES

\usepackage{graphicx}      %Voor het gebruik van figuren
\usepackage[utf8]{inputenc}%utf8 betekent dat naast de gewone 
                           %letters LaTeX ook letters zoals ë 
                           %á etc herkent %i.p.v.  {'o} voor ó

\usepackage{hyperref}     %Voor gebruik van links 

\usepackage{xcolor}       % voor gebruik van gekleurde tekst 
                          %probeer  \textcolor{red}{is dit rood?}
                          %Handig voor zichtbare herinneringen
\usepackage[dutch]{babel} %Voor het gebruik van Nederlandse taal
                          %en voor de automatische titels 
                          %zoals inhoudsopgave en referentielijst.

\usepackage{amsmath}      %Voor  wiskundige vergelijkingen
\usepackage{amssymb}      %Voor  wiskundige symbolen

\usepackage{tabularx}     %Voor het gebruik van tabellen
\usepackage{array}        %Voor multicol in tabellen

\usepackage{lipsum}   %voor lipsum dummytekst

%PACKAGES EN OPZET VOOR REFEREREN
\usepackage{csquotes}        % (nodig voor biblatex i.c.m. babel)
\usepackage[style=numeric, sorting=none]{biblatex}
\addbibresource{references.bib} %Zelfde naam als bib-bestand


%Specifiek toegevoegd voor het voorbeeld (staat dus niet in template)
\usepackage{pdfpages} %Voor het invoegen van PDF argumentatieschema

\title{Meeting 1}
\author{Tim Wiemer}
\date{2026/06/2026}

\begin{document}

\maketitle

\section{Checklist}

\begin{enumerate}
    \item{Was iedereen aanwezig} \\ Ja
    \item{Cijfer doelstelling} \\ Zo hoog mogelijk, 8
    \item{Sterke en zwakke punten}\\
    Tim: Coderen (S), Spelling en schrijven (Z),  \\
    Younes: Coderen (Z), Oplossend denken (S), Engels spreken (Z) \\ 
    Brecht: Coderen (Z), Formules herleiden (S)\\ 
    Amber: Coderen (Z), Creatief (S)\\
    \item{beschikbaarheid} elke dag van 9-17 beschikbaar, elke dag.\\
        Behalve Amber: donderdag en vrijdag helemaal niet.\\
        Tim: Woensdag 3juni enkel in de ochtend.
    \item{Rubric gelezen}, 
    \item{wat willen we leren}: 
        Wat is het leven van een onderzoeker, hoe ziet het dagelijksleven van een wetenschapper eruit. 
    \item{ Hebben we vragen voor de supervisor}:\\
        Wie bent u? \\
        Heeft u hobbies? \\
        Lievelingskleur? \\
        favoriete snack? \\
    \item{Vragen over de literatuur}\\
        Wat betekend de FWHM nou precies, hoe moeten we dat interpreteren?\\
        Wat als we nu wel gaan kijken naar de relativistische aspecten\\
    \item{Wat willen we bereiken}\\
        We willen meer practische ervaring opdoen. \\
        We willen leren hoe github werkt
    \item{Hebben we ervaring met Git / programeren}\\
        We hebben enkel ervaring met programmeren
    \item{Hoe was onze ervaring met DAS / NSP1}\\
        Das was prima en goed te doen. \\
        NSP1, veel geprogrameerd, best wel leuk. Labjournaal kostte veel moeite.
\end{enumerate}

\end{document}